% =============================================================
%  APPENDIX E — LIMITATIONS AND FUTURE DIRECTIONS
% =============================================================

\section{Limitations and Future Directions}
\label{app:limitations}

This appendix collects methodological caveats and directions for future 
work that emerge from the current implementation of RD$^*$\xspace.

% ------------------------------------------------------------------
\paragraph{Mean pooling vs.\ token-level analysis.}
% ------------------------------------------------------------------

The dimensionality and specialisation analyses of 
Appendix~\ref{app:head_analysis} rely on mean-pooled head representations 
$\mathbf{r}_{\ell,b,h}$ (Eq.~\eqref{eq:spatial_pooling}), collapsing 
intra-sample spatial variance to a single fixed-size vector per clip. 
This is motivated by computational tractability: constructing the full 
token-level matrix $\widetilde{\mathbf{R}}_{\ell,b,h} \in 
\mathbb{R}^{n\,N_w^\ell M \times d_h}$ for all $H_{\mathrm{tot}} = 184$ 
heads simultaneously would exceed available memory. 
The spectral reweighting operates at full spatial resolution, so this 
approximation affects only the interpretability analysis. 
Future work could address it by fitting the PCA basis on a token-level 
subsample.

% ------------------------------------------------------------------
\paragraph{Hook placement: pre- vs.\ post-projection.}
% ------------------------------------------------------------------

The current implementation intercepts raw head outputs 
$\mathbf{H}_{\ell,b,h}$ before the output projection $W^O_{\ell,b}$ 
(cf.\ \S\ref{app:arch_injection}). A complementary analysis of the 
post-projection contributions $\widehat{\mathbf{H}}_{\ell,b,h}$ yielded 
no meaningful differences in geometric structure, motivating the 
pre-projection formulation for its cleaner cross-stage comparability. 
A systematic comparison of the two placements under the full RD$^*$\xspace 
optimisation pipeline remains an open empirical question.

% ------------------------------------------------------------------
\paragraph{Dataset-specific PCA bases.}
% ------------------------------------------------------------------

Unlike the original ResiDual~\cite{residual2024}, which transfers a 
single PCA basis estimated on ImageNet across tasks, RD$^*$\xspace estimates 
$\{\boldsymbol{\Phi}_{\ell,b,h},\, \boldsymbol{\mu}_{\ell,b,h}\}$ from 
a small task-specific partition ($N_{\mathrm{PCA}} = 250$ samples per 
fold). Although the original paper shows dataset-specific and 
ImageNet-based bases to be largely equivalent in isotropic vision 
transformers (Table~8 of~\cite{residual2024}), we favour the 
task-specific strategy for two reasons. First, it avoids potential 
distribution mismatch: the acoustic heterogeneity of our benchmarks 
(environmental sounds, orchestral timbres, vocal categories) may not be 
well covered by any single general-purpose audio corpus. Second, it 
constitutes a more direct and targeted form of lightweight adaptation, 
analogous in spirit to fine-tuning but at a fraction of the 
computational cost.

% ------------------------------------------------------------------
\paragraph{Scope of the reweighting intervention.}
% ------------------------------------------------------------------

ResiDual~\cite{residual2024} applies spectral reweighting to all 
residual units, including MLP layers and, in some configurations, the 
output encoding ($\mathrm{RD}_Y$). RD$^*$\xspace intervenes exclusively on 
attention heads, reflecting the primary objective of this work: to 
characterise the internal geometry of HTS-AT heads and leverage it for 
targeted adaptation. Head-level intervention acts directly on the 
components analysed in Appendix~\ref{app:head_analysis}, avoiding the 
more entangled representations of MLP layers. Early experiments 
considered reweighting entire blocks and MLP sub-layers, but head-level 
intervention was preferred for its greater granularity and its direct 
correspondence with the geometric analysis. Extending RD$^*$\xspace to other 
residual components remains a natural direction for future work.

% ------------------------------------------------------------------
\paragraph{Incomplete stage-selection search.}
% ------------------------------------------------------------------

Table~\ref{tab:kfold_summary} evaluates only three configurations 
($\mathcal{L} \in \{\{3\},\{2,3\},\{0,1,2,3\}\}$); combinations such 
as $\{1,2\}$ or $\{1,2,3\}$ remain untested, and no formal statistical 
testing has been conducted. The reported results should therefore be 
interpreted as indicative rather than conclusive. A full combinatorial 
search with appropriate significance tests would be needed to establish 
the optimality of $\mathcal{L}=\{2,3\}$ with statistical rigour.

% ------------------------------------------------------------------
\paragraph{Extending the specialisation analysis.}
% ------------------------------------------------------------------

The cross-modal grounding of \S\ref{sec:specialization_method} 
identifies each head's dominant semantic axis via its first principal 
component alone. Jointly considering the top-$k$ components would 
provide a richer characterisation of heads encoding multiple overlapping 
concepts. Additionally, the qualitatively noted correlation between 
intrinsic dimensionality and semantic coherence Z-score has not been 
formally quantified; doing so would strengthen the theoretical grounding 
of the stage-selection criterion.

% ------------------------------------------------------------------
\paragraph{Computational overhead.}
% ------------------------------------------------------------------

The permanent forward hooks and per-block QKV re-computation 
(cf.\ footnote in \S\ref{app:arch_injection}) introduce a modest 
inference overhead at target blocks. A detailed latency profiling 
relative to the frozen CLAP baseline has not been included and would 
be a useful addition for deployment-oriented applications.